\scanpage{18}
\mspage{17}

On va préciser comment $l_g$ varie avec $g$. Soit $g'$ une autre section. \unsure{Alors|Nous} on a

\[
L(\lambda,g')=L(\lambda,g)\,g'{}^{*}(L(\lambda,g))^{-1}=
\]

\[
\text{\deleted{donc}}\qquad W(g')=W(g)\,g'_B{}^*(W(g))^{-1}
\]

(où $g'_B$ est la section de $X\times_S B/B$ déduite de la section $g'$ de $X/S$). On a donc

\[
\operatorname{pr}_1^{*}(L(\lambda,g'))W(g')=\operatorname{pr}_1^{*}(L(\lambda,g'))W(g)\,D
\]

\margin{$\begin{aligned}D={}&g'_B{}^{*}(L(\lambda,g))^{-1}\otimes_{\unsure{\mathcal O_X|O_X}}\\ &g'_B{}^*(W(g))^{-1}\end{aligned}$}

d'où comme d'habitude

\[
\begin{aligned}l_{g'}(\lambda)&=l_g(\lambda)+\bigl[\text{classe dans }\mathrm{Pic}_{B/S}\text{ de }g'_B{}^{*}(W(g))^{-\chi(\lambda)}\\ &\hspace{12em}=(N_{g'}N_g^{-1})^{-\chi(\lambda)}\\ &=l_g(\lambda)+(\psi_X(g')-\psi_X(g))(-\chi(\lambda))\end{aligned}\tag*{(*)}
\]

où $\psi_X\colon X\longrightarrow\underline{\mathrm{Alb}}^1_{X/S}$ est l'hom canonique, et $\underline{\mathrm{Alb}}^1_{X/S}$ est considéré comme torseur sous $A=\underline{\mathrm{Pic}}^0_{B/S}$ de la façon habituelle. Donc, considérant l'hom

\[
\varphi\colon\underline{\mathrm{NS}}_{X/S}\longrightarrow\underline{\mathrm{NS}}_{B/S}
\]

induit par $l$ de façon évidente, on déduit pour tte section $\delta$ de $\underline{\mathrm{NS}}_{X/S}$ un hom

\[
\boxed{l_\delta\colon X\times\underline{\mathrm{Pic}}^\delta_{X/S}\longrightarrow\underline{\mathrm{Pic}}^{\varphi(\delta)}_{B/S}}
\]

\deleted{qui, \unsure{sur} définit \unsure{; se donne}} \deleted{se factorise $\underline{\mathrm{Alb}}^{\deleted{\unsure{-\chi(\delta)|\chi(\delta)}}}_{X/S}$ d'ailleurs par $\underline{\mathrm{Alb}}^{-\chi(\delta)}_{X/S}\overset{A}{\times}$} satisfaisant \unsure{aux,|à,|les,}
